\chapter{Benchmarks and latency appendix}\label{benchmarks-and-latency-appendix}

\section{Section summary}\label{section-summary}

Observed numbers frame how the stack performs. The headline comparison is the 41-turn multi-agent bench (12 October 2026 → 14 October 2026): finish-reason mismatches fell from 6/41 to 0/41, retrieval hit rose from 0.537 to 0.780, and p50 latency fell 39\% from 8006 ms to 4872 ms. Published knowledge-base quality on the frozen v11 export is \textbf{Hit@5 0.9306}. The current corpus is v13 (\texttt{1405-07-06}), and warm hybrid search stays near \textbf{55 ms}. End-to-end chat latency is led by guarded generation on Gemma.

\section{Knowledge-base IR quality}\label{knowledge-base-ir-quality}

\textbf{Sources:} KB pages corpus table; \texttt{components/knowledgebase/kb-manager/artifacts/retrieval\_training/report\_metrics.json} (\texttt{hit5\_rate} 0.9306, \texttt{mrr} 0.8281, \texttt{N} 663). The pages also describe a frozen 120-query, six-format set that carries the same v11 Hit@5. The checked-in rate is the 663-query export. v13 (\texttt{1405-07-06}, 18 documents / 1,123 chunks) is the RAM-first corpus whose warm latency is reported below.

\begin{longtable}[]{@{}
  >{\raggedright\arraybackslash}p{(\columnwidth - 6\tabcolsep) * \real{0.2143}}
  >{\raggedright\arraybackslash}p{(\columnwidth - 6\tabcolsep) * \real{0.2143}}
  >{\raggedleft\arraybackslash}p{(\columnwidth - 6\tabcolsep) * \real{0.2857}}
  >{\raggedleft\arraybackslash}p{(\columnwidth - 6\tabcolsep) * \real{0.2857}}@{}}
\toprule\noalign{}
Snapshot & Corpus & \begin{minipage}[b]{\linewidth}\raggedleft
Hit@5
\end{minipage} & \begin{minipage}[b]{\linewidth}\raggedleft
MRR
\end{minipage} \\
\midrule\noalign{}
\endhead
\bottomrule\noalign{}
\endlastfoot
v11 & 35 docs / 2,227 chunks; \texttt{dense\ use\_context=False} & \textbf{0.9306} & 0.8281 \\
\end{longtable}

\section{Knowledge-base stage latency (ms)}\label{knowledge-base-stage-latency-ms}

\textbf{Source:} \texttt{artifacts/retrieval\_training/latency\_v13\_1.json}

\begin{longtable}[]{@{}
  >{\raggedright\arraybackslash}p{(\columnwidth - 6\tabcolsep) * \real{0.2000}}
  >{\raggedleft\arraybackslash}p{(\columnwidth - 6\tabcolsep) * \real{0.2667}}
  >{\raggedleft\arraybackslash}p{(\columnwidth - 6\tabcolsep) * \real{0.2667}}
  >{\raggedleft\arraybackslash}p{(\columnwidth - 6\tabcolsep) * \real{0.2667}}@{}}
\toprule\noalign{}
Stage & \begin{minipage}[b]{\linewidth}\raggedleft
README v13 (10-07)
\end{minipage} & \begin{minipage}[b]{\linewidth}\raggedleft
Before tune (10-13)
\end{minipage} & \begin{minipage}[b]{\linewidth}\raggedleft
After tune (10-13)
\end{minipage} \\
\midrule\noalign{}
\endhead
\bottomrule\noalign{}
\endlastfoot
BM25 & 4.0 & 3.8 & 2.6 \\
Dense & 24.0 & 20.5 & 13.5 \\
Rerank & 24.0 & 20.9 & 22.0 \\
\textbf{Total} & \textbf{55.0} & \textbf{44.4} & \textbf{39.6} \\
\end{longtable}

The v13 README figure of \textbf{55.0 ms} is the warm-search number to quote for the RAM-first stack (\textasciitilde55 ms). The 13 October tune later measured 39.6 ms p50 on 30 fresh queries; both describe retrieval, not chat. Cold starts remain multi-second when models and indexes load. Redis cache hits for identical queries fall to the low single-digit milliseconds once warm.

\begin{figure}
\centering
\includegraphics[width=0.88\linewidth,height=0.36\textheight,keepaspectratio]{../assets/plots/latency_breakdown_ram_tantivy.png}
\caption{KB latency RAM/tantivy plot}
\end{figure}

\textbf{Figure B1. KB latency breakdown (RAM / tantivy era).} The plot visualizes how lexical and dense legs plus rerank compose wall time after memory-resident indexes replaced slower disk-bound paths. Read it with the table above: the README warm total is about 55 ms, and the 13 October tune sits under 50 ms. Either figure is retrieval time, not chat generation.

\section{End-to-end orchestrator bench (41 turns)}\label{end-to-end-orchestrator-bench-41-turns}

\textbf{Source:} parent pages \texttt{docs/papers/benchmarks.html} (before 2026-10-12 → after 2026-10-14). Reproduce with \texttt{evaluation.bench\_runner} → \texttt{evaluation.ragas\_eval} → \texttt{evaluation.report}, which writes \texttt{bench\_report\_latest}. The after column's overall p50 (4872 ms) matches that artifact's headline; the before column and the per-metric deltas below are the published comparison, not a re-fit.

\begin{longtable}[]{@{}lrrl@{}}
\toprule\noalign{}
Metric & Before (2026-10-12) & After (2026-10-14) & Δ \\
\midrule\noalign{}
\endhead
\bottomrule\noalign{}
\endlastfoot
Finish mismatches & 6/41 & 0/41 & −6 \\
Retrieval hit & 0.537 & 0.780 & +0.243 \\
Cheque hit & 0.111 & 0.667 & +0.556 \\
Context recall & 0.113 & 0.340 & +0.227 \\
Answer relevancy & 0.212 & 0.568 & +0.356 \\
Ambiguity resolution & 0.854 & 1.000 & +0.146 \\
Sticky persistence & 0.610 & 0.683 & +0.073 \\
p50 latency & 8006 ms & 4872 ms & \textbf{−39\%} \\
\end{longtable}

The same after cut, by subspace:

\begin{longtable}[]{@{}lrrrrr@{}}
\toprule\noalign{}
Subspace & n & Hit & Faith & Sticky & p50 \\
\midrule\noalign{}
\endhead
\bottomrule\noalign{}
\endlastfoot
Cheque & 9 & 0.667 & 0.422 & 0.444 & 10301 ms \\
Facilities person & 6 & 0.500 & 0.967 & 0.667 & 10718 ms \\
Facilities legal & 3 & 0.667 & 0.917 & 0.333 & 10498 ms \\
General & 10 & 0.800 & 0.655 & 0.600 & 6981 ms \\
\end{longtable}

\texttt{bench\_report\_latest} (41 turns, 12 conversations, model \texttt{ics-helper-agent}, 2026-10-14) adds figures the delta table does not carry: average latency \textbf{4938 ms}, transport error rate \textbf{0\%}, faithfulness \textbf{0.461}, context relevancy \textbf{0.283}, hallucination inverted \textbf{0.746}. Where a per-subspace p50 in an earlier export disagrees with the table above, the benchmarks page is the comparison this report uses.

Generation through guardrails accounts for most of the multi-second p50; warm KB search (\textasciitilde55 ms, and 39.6 ms after the 13 October tune) is two orders of magnitude smaller. The −39\% p50 move (8006 → 4872) places chat near five seconds. Throughput work centers on TRT-LLM on the H200, prompt and entity trimming, and the clarify path, beside the warm BM25 leg.

\section{Model serving concurrency (H200)}\label{model-serving-concurrency-h200}

\textbf{Source:} \texttt{components/server-setup/docs/history/013\_2026-10-13\_trtllm\_champion\_concurrency.md} (as summarized in ops notes)

\begin{longtable}[]{@{}
  >{\raggedright\arraybackslash}p{(\columnwidth - 4\tabcolsep) * \real{0.3333}}
  >{\raggedright\arraybackslash}p{(\columnwidth - 4\tabcolsep) * \real{0.3333}}
  >{\raggedright\arraybackslash}p{(\columnwidth - 4\tabcolsep) * \real{0.3333}}@{}}
\toprule\noalign{}
Backend & Approx. aggregate throughput & Notes \\
\midrule\noalign{}
\endhead
\bottomrule\noalign{}
\endlastfoot
llama.cpp Q4 Gemma pool & \textasciitilde58 tok/s class & Multi-process supervisor \\
TRT-LLM BF16 champion & \textasciitilde3874 tok/s @ concurrency 128 & Large aggregate gain vs llama.cpp \\
\end{longtable}

Exact reproduction depends on engine build and GPU partitioning. The ratio (\textasciitilde67× aggregate in that experiment) is the published evidence that \textbf{scale-up on the H200 prefers TRT-LLM for generation}, with the llama.cpp pool on \texttt{8080–8084} as the resilience path.
